\section*{Problemas}

\subsection*{Enunciados de los problemas}

% Recordar
\begin{problema}
	\label{prob:de3bea1}
	Clasifique cada una de las siguientes variables según su escala de medición (\emph{Nominal, Ordinal, de Intervalo o de Razón}):
	\begin{enumerate}
		\item El código postal de la residencia de un cliente en una plataforma de comercio electrónico.
		\item El nivel de satisfacción de un usuario respecto a una aplicación móvil (Muy Insatisfecho, Insatisfecho, Neutral, Satisfecho, Muy Satisfecho).
		\item La temperatura en grados Celsius (°C) registrada por los servidores de un centro de datos.
		\item El salario mensual en pesos mexicanos de un ingeniero en ciencia de datos.
	\end{enumerate}
\end{problema}

% Comprender
\begin{problema}
	\label{prob:7b7dc88}
	Determine si la métrica descrita en los siguientes enunciados representa un \textbf{Parámetro} ($\mu, \sigma^2, p$) o un \textbf{Estadístico} ($\bar{X}, S^2, \hat{p}$), y explique el criterio que distingue ambos casos:
	\begin{enumerate}
		\item Tras censar las 32 entidades federativas, el INEGI determinó que el promedio de personas por vivienda en el país es de $3.6$.
		\item Un ingeniero de calidad selecciona un lote aleatorio de 80 procesadores y calcula una varianza en la frecuencia de reloj de $0.015 \text{ GHz}^2$.
		\item En un subconjunto de prueba de 5,000 transacciones bancarias, la tasa de fraude detectada fue del $0.84\%$.
	\end{enumerate}
\end{problema}

% Aplicar
\begin{problema}
	\label{prob:4bc23ad}
	Una corporación de consultoría tecnológica cuenta con 4,000 empleados divididos en cuatro divisiones: Desarrollo ($2,000$), Infraestructura ($1,000$), Ventas ($600$) y Recursos Humanos ($400$). Se requiere seleccionar una muestra aleatoria de $n = 200$ empleados para evaluar el clima laboral. Diseñe un esquema de \emph{muestreo aleatorio estratificado proporcional} y determine exactamente cuántos empleados deben entrevistarse en cada una de las cuatro divisiones.
\end{problema}

% Analizar
\begin{problema}
	\label{prob:a5bddfb}
	En 1936, la revista \emph{Literary Digest} realizó un muestreo masivo enviando 10 millones de cuestionarios por correo para predecir la elección presidencial de Estados Unidos, utilizando como marco muestral directorios telefónicos y registros de propietarios de automóviles. A pesar de obtener una muestra gigantesca de $n \approx 2.4$ millones de respuestas, predijeron un triunfo contundente de Alf Landon, cuando en realidad Roosevelt ganó por avalancha. Demuestre conceptual y algebraicamente, mediante la descomposición $\text{MSE}(\bar{X}_n) = \sigma^2/n + B^2$, cómo un sesgo sistemático $B = \mathbb{E}[\bar{X}] - \mu \neq 0$ hace que aumentar $n \to \infty$ no reduzca el error total, sino que la estimación converja con certeza hacia un valor incorrecto.
\end{problema}

% Evaluar
\begin{problema}
	\label{prob:feeec4c}
	Para evaluar el éxito de una nueva plataforma de banca web, un departamento de experiencia de usuario manda un correo masivo a 100,000 clientes invitándolos a contestar voluntariamente un cuestionario en línea sobre usabilidad. Responden 1,200 clientes, y el $88\%$ de ellos manifiesta que la plataforma es extremadamente confusa. Evalúe qué sesgo metodológico predomina en esta recolección de datos y por qué distorsiona la proporción real del total de clientes ($p$), y proponga un rediseño del protocolo de muestreo para obtener una estimación insesgada sin incrementar el presupuesto de recolección.
\end{problema}

% Crear
\begin{problema}
	\label{prob:701b535}
	Diseñe un ejemplo original (distinto de los casos de \emph{Literary Digest}, streaming o banca web ya vistos) de un estudio de ciencia de datos donde el mecanismo de recolección introduzca un sesgo sistemático de cobertura o de autoselección. Identifique explícitamente la población objetivo, la muestra realmente obtenida, el tipo de sesgo presente, y proponga un rediseño del protocolo de muestreo que lo elimine o mitigue.
\end{problema}

\subsection*{Soluciones detalladas}

% Recordar
\begin{solproblema}[prob:de3bea1]
	La clasificación por nivel o escala de medición se justifica de la siguiente forma:
	\begin{enumerate}
		\item \textbf{Código postal $\to$ Nominal:} Los números actúan exclusivamente como etiquetas o categorías de ubicación geográfica. No existe un orden natural ni sentido matemático en sumar o restar códigos postales.
		\item \textbf{Satisfacción de usuario $\to$ Ordinal:} Existe una jerarquía o clasificación intrínseca (Muy Insatisfecho $<$ Insatisfecho $<$ Neutral $<$ Satisfecho $<$ Muy Satisfecho), pero la distancia conceptual entre una categoría y otra no es necesariamente uniforme ni cuantificable.
		\item \textbf{Temperatura en °C $\to$ De Intervalo:} La diferencia numérica de grados tiene un significado físico constante, pero el cero es arbitrario (el $0^\circ\text{C}$ representa el punto de congelación del agua, no la ausencia absoluta de energía térmica). Por tanto, no podemos afirmar que $40^\circ\text{C}$ sea el "doble de calor" que $20^\circ\text{C}$.
		\item \textbf{Salario mensual $\to$ De Razón:} Posee intervalos uniformes y un cero absoluto real ($0$ pesos implica carencia total de ingresos salariales). Las razones son válidas: un salario de $\$60,000$ es exactamente el doble de uno de $\$30,000$.
	\end{enumerate}
\end{solproblema}

% Comprender
\begin{solproblema}[prob:7b7dc88]
	La distinción entre parámetro y estadístico se basa en si la métrica describe al universo total censado o a una porción del mismo:
	\begin{enumerate}
		\item \textbf{Promedio de $3.6$ personas por vivienda $\to$ Parámetro ($\mu$):} Proviene de un censo nacional exhaustivo que abarcó la totalidad de las 32 entidades federativas.
		\item \textbf{Varianza de $0.015 \text{ GHz}^2$ $\to$ Estadístico ($S^2$):} Es una medida calculada exclusivamente sobre la muestra aleatoria de 80 procesadores inspeccionados.
		\item \textbf{Tasa de fraude del $0.84\%$ $\to$ Estadístico ($\hat{p}$):} Es una proporción estimada sobre el subconjunto muestral de 5,000 transacciones extraídas para prueba.
	\end{enumerate}
	El criterio decisivo no es la magnitud del conjunto de datos, sino si se originó de un censo exhaustivo (parámetro) o de una porción aleatoria del universo (estadístico).
\end{solproblema}

% Aplicar
\begin{solproblema}[prob:4bc23ad]
	En un muestreo aleatorio estratificado proporcional, el tamaño de muestra de cada estrato $k$ es proporcional a su peso en la nómina: $n_k = n \cdot \left(\frac{N_k}{N}\right) = 200 \cdot \left(\frac{N_k}{4,000}\right)$. Sustituimos:
	\begin{itemize}
		\item \textbf{Desarrollo:} $n_1 = 200 \cdot \left(\frac{2,000}{4,000}\right) = 100$ empleados.
		\item \textbf{Infraestructura:} $n_2 = 200 \cdot \left(\frac{1,000}{4,000}\right) = 50$ empleados.
		\item \textbf{Ventas:} $n_3 = 200 \cdot \left(\frac{600}{4,000}\right) = 30$ empleados.
		\item \textbf{Recursos Humanos:} $n_4 = 200 \cdot \left(\frac{400}{4,000}\right) = 20$ empleados.
	\end{itemize}
	La suma total verifica que $100 + 50 + 30 + 20 = 200$ entrevistas, con paridad estructural garantizada respecto al peso real de cada división.
\end{solproblema}

% Analizar
\begin{solproblema}[prob:a5bddfb]
	Sea $\bar{X}_n$ el estimador muestral basado en una muestra de tamaño $n$. El error cuadrático medio (MSE) de un estimador descompone la imprecisión total en varianza muestral y sesgo al cuadrado:
	\begin{align}
		\text{MSE}(\bar{X}_n) = \mathbb{E}\left[(\bar{X}_n - \mu)^2\right] = \text{Var}(\bar{X}_n) + (\text{Sesgo})^2 = \dfrac{\sigma^2}{n} + B^2,
	\end{align}
	donde $B = \mathbb{E}[\bar{X}_n] - \mu$ representa el sesgo metodológico derivado del marco de selección (en \emph{Literary Digest}, encuestar solo a personas pudientes con teléfono y automóvil en plena Gran Depresión).

	Al incrementar el tamaño de muestra indefinidamente ($n \to \infty$), el componente de variabilidad muestral tiende a cero ($\lim_{n \to \infty} \sigma^2/n = 0$), por la Ley de los Grandes Números. Sin embargo, el componente de sesgo sistemático permanece exactamente igual:
	\begin{align}
		\lim_{n \to \infty} \text{MSE}(\bar{X}_n) = 0 + B^2 = B^2 > 0.
	\end{align}
	Esto demuestra que una muestra gigante de millones de registros sesgados ($n \approx 2.4$ millones) únicamente reduce la varianza en torno a un valor incorrecto, generando una falsa ilusión de precisión sobre una estimación errónea. En contraste, una muestra pequeña como la de Gallup ($n = 50,000$), donde $B \approx 0$, alcanza una convergencia insesgada hacia el verdadero parámetro $\mu$.
\end{solproblema}

% Evaluar
\begin{solproblema}[prob:feeec4c]
	Predomina el \textbf{sesgo de autoselección o respuesta voluntaria}. Aquellos usuarios que experimentaron frustración técnica grave con la plataforma tienen una motivación psicológica mucho mayor para abrir la encuesta y quejarse que aquellos usuarios satisfechos o indiferentes que completaron sus operaciones sin problema. Por tanto, el $88\%$ es una sobrestimación masiva del nivel real de confusión poblacional ($p$).

	\textbf{Rediseño sin costo extra:} en lugar de mandar un correo pasivo a 100,000 personas, se debe seleccionar una muestra aleatoria pura mucho más pequeña (por ejemplo, $n = 500$ clientes) y aplicar un seguimiento activo intermitente en la misma sesión web (\emph{pop-up} aleatorio controlado al finalizar una transacción) o contacto directo programado, asegurando una tasa de respuesta alta y no sesgada por frustración personal.
\end{solproblema}

% Crear
\begin{solproblema}[prob:701b535]
	\textbf{Ejemplo:} una universidad quiere estimar el promedio de horas semanales que sus egresados dedican a educación continua, y para ello envía la encuesta únicamente a través de su grupo oficial de LinkedIn. \textbf{Población objetivo:} todos los egresados de la universidad. \textbf{Muestra obtenida:} solo los egresados que mantienen una cuenta activa de LinkedIn y siguen la página institucional — un subconjunto sesgado hacia quienes están profesionalmente más activos y conectados. \textbf{Tipo de sesgo:} sesgo de cobertura (el marco muestral —usuarios de LinkedIn— excluye sistemáticamente a egresados sin presencia en esa red, típicamente quienes cambiaron de carrera o se desvincularon del mercado laboral formal), agravado por un posible sesgo de autoselección entre quienes sí ven la publicación. \textbf{Rediseño:} usar el directorio institucional completo de egresados (correo electrónico registrado al titularse) como marco muestral, y extraer de ahí una muestra aleatoria simple, en lugar de depender de la membresía voluntaria a una red social.
\end{solproblema}
